% =========================================================================
% PLANTILLA MAESTRA: TRATADO ACADÉMICO FORMAL / TESIS / MONOGRAFÍA
% =========================================================================
% Optimizada para monografías formales, informes de investigación, tesis
% de grado y publicaciones académicas tradicionales.
%
% Características Editoriales:
% 1. Portada formal institucional con metadatos (\maketitle).
% 2. Resumen ejecutivo estructurado (\begin{abstract}).
% 3. Tabla de contenidos interactiva (\tableofcontents).
% 4. Geometría estándar ISO A4 con márgenes de 2.5 cm.
% 5. Tipografía clara a 11pt con soporte microtipográfico.
% 6. Tablas limpias con booktabs y soporte matemático robusto (amsmath/amssymb).
% =========================================================================

\documentclass[11pt,a4paper]{article}

% --- Codificación y Lenguaje ---
\usepackage[utf8]{inputenc}
\usepackage[T1]{fontenc}
\usepackage[spanish,es-tabla,es-noquoting]{babel}

% --- Geometría y Márgenes ---
\usepackage{geometry}
\geometry{
    top=2.5cm,
    bottom=2.5cm,
    left=2.5cm,
    right=2.5cm
}

% --- Matemáticas y Símbolos ---
\usepackage{amsmath,amssymb,amsfonts}

% --- Gráficos y Tablas ---
\usepackage{graphicx}
\usepackage{booktabs}
\usepackage{tabularx}
\usepackage{xcolor}

% --- Microtipografía y Legibilidad ---
\usepackage{microtype}

% --- Enlaces y Metadatos del PDF ---
\usepackage{hyperref}
\hypersetup{
    colorlinks=true,
    linkcolor=blue!80!black,
    urlcolor=blue!80!black,
    citecolor=blue!80!black,
    pdfauthor={Hugo Joaquín Hidalgo},
    pdftitle={Tratado Académico Formal},
    pdfsubject={Monografía / Tesis / Informe de Investigación},
    pdfcreator={pdflatex}
}

% --- Metadatos del Documento ---
\title{\vspace{-1cm}\textbf{\LARGE Tratado Académico Formal}\\[0.5em] \large Subtítulo Descriptivo o Alcance del Trabajo}
\author{\textbf{Hugo Joaquín Hidalgo}\\[0.2em]
\small Facultad de Ingeniería --- Universidad Nacional de Jujuy\\
\small \texttt{contacto@ejemplo.com}}
\date{\today}

% =========================================================================
% CUERPO DEL DOCUMENTO
% =========================================================================
\begin{document}

\maketitle

\begin{abstract}
Este documento ejemplifica la estructura formal requerida para monografías académicas, informes técnicos de investigación y proyectos finales de grado. Se presenta un resumen ejecutivo conciso del contexto problemático, los objetivos metodológicos, los resultados analíticos cuantitativos y las conclusiones principales alcanzadas, permitiendo al evaluador una rápida comprensión de la contribución técnica.
\end{abstract}

\vspace{1em}
\noindent\textbf{Palabras clave:} Computación de Alto Rendimiento, Redes de Interconexión, Latencia, Arquitectura de Sistemas, Rendimiento.

\newpage
\tableofcontents
\newpage

% --- SECCIÓN 1: INTRODUCCIÓN ---
\section{Introducción}
La creciente demanda de procesamiento en sistemas de cómputo modernos exige arquitecturas capaces de mitigar el cuello de botella de la memoria mediante topologías de comunicación escalables y jerarquías de almacenamiento optimizadas.

\subsection{Motivación y Contexto}
El modelado analítico de redes de interconexión en multiprocesadores de chip (CMP) permite predecir el comportamiento del sistema ante saturación de tráfico, optimizando la asignación de buffers y el consumo energético de los enrutadores.

\subsection{Objetivos del Trabajo}
\begin{itemize}
    \item Formalizar el modelo matemático de latencia de red en topologías de malla 2D y torus.
    \item Evaluar el impacto de la contención de buffers mediante análisis estocástico.
    \item Sintetizar recomendaciones directas de diseño para implementaciones en silicio.
\end{itemize}

% --- SECCIÓN 2: MARCO TEÓRICO ---
\section{Marco Teórico y Fundamentos Matemáticos}
La latencia base de transporte $T_{\text{base}}$ de un paquete de $L$ flits en una red con $H$ saltos promedio se modela formalmente como:

\begin{equation}
    T_{\text{base}} = H \cdot t_r + \frac{L}{W} + H \cdot t_w
\end{equation}

\noindent donde:
\begin{itemize}
    \item $H$: Número medio de saltos de la topología (\textit{hop count}).
    \item $t_r$: Latencia interna del pipeline del router (ciclos de reloj).
    \item $W$: Ancho de canal por enlace (bits/ciclo).
    \item $t_w$: Retardo de propagación por el cable físico entre nodos adyacentes.
\end{itemize}

% --- SECCIÓN 3: RESULTADOS Y COMPARATIVA ---
\section{Evaluación y Resultados Experimentales}
La Tabla~\ref{tab:topologias} resume las métricas topológicas fundamentales comparando configuraciones estándar para $N = 64$ nodos.

\begin{table}[htbp]
\centering
\caption{Comparativa analítica de topologías de interconexión ($N=64$).}
\label{tab:topologias}
\begin{tabularx}{\textwidth}{Xcccc}
\toprule
\textbf{Topología} & \textbf{Grado ($d$)} & \textbf{Diámetro ($D$)} & \textbf{Bisección ($B$)} & \textbf{Complejidad Cableado} \\
\midrule
Malla 2D ($8 \times 8$) & 4 & 14 & 8 & Baja (Planar) \\
Torus 2D ($8 \times 8$) & 4 & 8  & 16 & Media (Enlaces envolventes) \\
Hipercubo ($n=6$)      & 6 & 6  & 32 & Alta (Logarítmica) \\
Crossbar Centralizado   & 64 & 1 & 64 & Muy Alta ($\mathcal{O}(N^2)$) \\
\bottomrule
\end{tabularx}
\end{table}

% --- SECCIÓN 4: CONCLUSIONES ---
\section{Conclusiones}
La topología de malla 2D ofrece el balance óptimo entre complejidad de enrutamiento físico sobre el die y rendimiento para tamaños de clúster moderados. Para escalas mayores a 64 nodos, el torus doblado (\textit{folded torus}) reduce el diámetro a la mitad sin introducir enlaces largos asimétricos.

% --- REFERENCIAS BIBLIOGRÁFICAS ---
\begin{thebibliography}{99}
\bibitem{dally2004} Dally, W. J., \& Towles, B. (2004). \textit{Principles and Practices of Interconnection Networks}. Morgan Kaufmann Publishers.
\bibitem{hennessy2019} Hennessy, J. L., \& Patterson, D. A. (2019). \textit{Computer Architecture: A Quantitative Approach} (6th ed.). Elsevier.
\end{thebibliography}

\end{document}
